% Capítulo textual do trabalho.

\chapter{Cronograma}
\label{cap:cronograma}

O cronograma é uma ferramenta essencial no planejamento de um Trabalho de Conclusão de Curso~(TCC). Ele representa a programação das atividades necessárias para a conclusão do projeto, facilitando a organização e o gerenciamento do tempo. Um cronograma bem elaborado é fundamental para garantir que todas as etapas do TCC sejam realizadas dentro do prazo, evitando sobrecargas de trabalho.

O cronograma deve incluir as principais etapas do TCC, como a definição do tema, revisão bibliográfica, desenvolvimento da metodologia, coleta e análise de dados, redação do trabalho e entrega final. O Quadro~\ref{tab:cronograma} apresenta um exemplo de cronograma.

\begin{quadro}[h]
    \legenda[tab:cronograma]{Cronograma detalhado das atividades do TCC}
    \begin{tabular}{|p{7cm}|c|c|c|c|c|}
        \hline
        \textbf{Atividade}                      & \textbf{Ago} & \textbf{Set} & \textbf{Out} & \textbf{Nov} & \textbf{Dez} \\
        \hline\hline
        Definição do Tema                      & \cellcolor{black}  &  &  &  &  \\\hline
        Revisão Bibliográfica                  & \cellcolor{black}  & \cellcolor{black}  &  &  &  \\\hline
        Desenvolvimento da Metodologia         &  & \cellcolor{black}  & \cellcolor{black}  &  &  \\\hline
        Coleta de Dados                        &  &  & \cellcolor{black}  & \cellcolor{black}  &  \\\hline
        Análise dos Dados                      &  &  &  & \cellcolor{black}  & \cellcolor{black}  \\\hline
        Redação do TCC                        &  &  &  & \cellcolor{black}  & \cellcolor{black}  \\\hline
        Revisão Final e Entrega                &  &  &  &  & \cellcolor{black}  \\
        \hline
    \end{tabular}
    \vspace{0.2cm}
    \fonteautor
\end{quadro}
